%%%%%%%%%%%%%%%%%%%%%%%%%%%%%%%%%%%%%%%%%%%%%%%%%%%%%%%%%%%%%%%%%%%%%%%%
% PÁGINAS PRELIMINARES
% Preámbulo, agradecimientos, dedicatoria, resumen y abstract.
%%%%%%%%%%%%%%%%%%%%%%%%%%%%%%%%%%%%%%%%%%%%%%%%%%%%%%%%%%%%%%%%%%%%%%%%

%% =====================================================================
%% PREÁMBULO (Justificación y Objetivos) -- obligatorio
%% =====================================================================
\chapter*{Preámbulo}
\addcontentsline{toc}{chapter}{Preámbulo}

La vigilancia de espacios interiores cerrados (oficinas, naves, almacenes o
viviendas) recae habitualmente en cámaras fijas y en la supervisión humana, dos
recursos con cobertura limitada y elevado coste de atención continuada. La
robótica móvil de servicio, unida a los avances recientes en visión por
computador, abre la posibilidad de automatizar esta tarea mediante un agente
capaz de recorrer el entorno de forma autónoma y de avisar únicamente cuando
detecta una situación relevante. El presente Trabajo Fin de Grado se enmarca en
este contexto y propone el diseño y la implementación de un sistema robótico de
patrulla autónoma para interiores, construido sobre un robot \mbox{TurtleBot~4} y
\gls{ros}, que cartografía el espacio, lo recorre siguiendo una ruta de
cobertura, detecta la presencia de personas mediante una red neuronal de
detección de objetos y notifica las incidencias a una plataforma centralizada
con la que el usuario interactúa desde una aplicación móvil. El objetivo general
es demostrar la viabilidad de una arquitectura desacoplada y orientada a eventos
que integre la capa robótica, el procesamiento a bordo del robot y los servicios web
en un único sistema funcional, extensible y seguro.

\cleardoublepage

%% =====================================================================
%% AGRADECIMIENTOS (opcional)
%% =====================================================================
\chapter*{Agradecimientos}
\addcontentsline{toc}{chapter}{Agradecimientos}

En primer lugar, quiero agradecer a mi tutor, Francisco Antonio Pujol López, su
orientación y su disponibilidad durante el desarrollo de este trabajo.

A mis padres y a mi pareja María, por confiar siempre en mí, incluso cuando yo no lo
hacía.


A mis amigos Alejandro, Miguel, Nicolás, Juan, Raúl Luján y Raúl Ruiz, por el apoyo
constante a lo largo de estos años y ser un pilar importantísimo en mi historia
como ingeniero.

Finalmente, agradecer a la \EPSfacultad{} de la \EPSuniversidad{} la inestimable
formación recibida durante el grado.

\cleardoublepage

%% =====================================================================
%% DEDICATORIA (opcional) -- una hoja a la derecha, sin entrada en índice.
%% Personaliza el texto; es un marcador de posición.
%% =====================================================================
\thispagestyle{empty}
\null\vfill
\begin{flushright}
  \itshape
  A mi familia y pareja,\\
  por su apoyo incondicional y por mantenerme siempre con los pies en la tierra.
\end{flushright}
\vfill
\cleardoublepage

%% =====================================================================
%% CITA / FRASE CÉLEBRE (opcional) -- una hoja a la derecha, sin índice.
%% Sustituye la cita y su autor por los que prefieras.
%% =====================================================================
\thispagestyle{empty}
\null\vfill
\begin{flushright}
  \itshape
  ``Solo podemos ver poco del futuro,\\ 
  pero lo suficiente para darnos cuenta de que hay mucho por hacer.''
  \par\bigskip
  \normalfont\small Alan Turing, 1950
\end{flushright}
\vfill
\cleardoublepage

%% =====================================================================
%% RESUMEN
%% =====================================================================
\chapter*{Resumen}
\addcontentsline{toc}{chapter}{Resumen}

La vigilancia de espacios interiores cerrados depende en gran medida de cámaras
fijas y de supervisión humana, soluciones de cobertura limitada y atención
costosa. Este trabajo aborda la automatización de dicha tarea mediante un robot
móvil capaz de patrullar de forma autónoma y de alertar solo ante eventos
relevantes. Para ello se ha diseñado e implementado un sistema completo sobre un
robot \mbox{TurtleBot~4} y \gls{ros}, organizado como un monorepo con cuatro
bloques: la capa robótica, una infraestructura de datos en contenedores, un
servicio de backend y una aplicación de usuario. El robot construye el mapa del
entorno mediante \gls{slam}, genera una ruta de cobertura a partir de dicho mapa
y la recorre con la pila de navegación \gls{nav2}, mientras un nodo de percepción
basado en \gls{yolo} detecta personas y registra las incidencias. La
comunicación entre el robot y el servidor se resuelve con una arquitectura
orientada a eventos basada en un bus de mensajería, lo que desacopla el hardware
del plano de control y evita la dependencia de una conexión permanente entre ambos
extremos. El backend, desarrollado con
FastAPI, consume estos eventos, persiste los datos y expone una \gls{api}
\gls{rest} consumida por una aplicación \mbox{Flutter} desde la que el usuario
controla la patrulla y revisa los incidentes. Los resultados presentados son de
carácter preliminar y se centran en la validación funcional de la arquitectura;
no obstante, el sistema sienta una base sólida y extensible para el desarrollo de
soluciones de vigilancia robótica autónoma.

\textbf{Palabras clave:} robótica móvil, ROS 2, navegación autónoma, visión por
computador, arquitectura orientada a eventos, vigilancia.

\cleardoublepage
